\guidelinesubsubsection{Recommendations}

Benchmarks, baselines, and metrics play an important role in assessing the effectiveness of LLMs and LLM-based tools. \added{A \textit{benchmark} is a ``standard tool for the competitive evaluation and comparison of competing systems or components according to specific characteristics, such as performance, dependability, or security''~\cite{Kistowski2015Benchmark} }
Benchmarks can be used to assess LLM performance on specific tasks such as code summarization, code generation, \added{and static analysis}.
Meanwhile, ``a \textit{metric} is a method, algorithm, or procedure for assigning one or more numbers to a phenomenon''~\cite{ralph2024teaching}. 
A baseline represents a reference point for the measured LLM.
Since LLMs require substantial hardware resources, baselines allow research to compare LLMs against traditional algorithms with lower computational costs.

When selecting benchmarks and metrics, it is important to fully understand the benchmark tasks, what exactly is being measured, and how it relates to the (often latent) variables researchers actually care about. If one or more metrics or benchmarks are used, researchers:
\begin{itemize}
    \item \must briefly justify (in the \paper) why selected benchmarks and metrics are suitable for the given task or study. 
    % \item \must discuss the reliability and validity (especially construct validity) of selected metrics and benchmarks
    \item \should summarize the structure and tasks of the selected benchmark(s), including the programming language(s) and descriptive statistics such as the number of contained tasks and test cases;
    \item \should discuss the limitations of the selected benchmark(s) (e.g. widely used benchmarks such as \emph{HumanEval} \cite{DBLP:journals/corr/abs-2107-03374} and \emph{MBPP} \cite{DBLP:journals/corr/abs-2108-07732} Python functions assess a very specific part of software development, which is not representative of the full breadth of SE work~\cite{Chandra2025benchmarks}).
    \item \added{\should} include an example of a task and the corresponding test case(s) to illustrate the structure of the benchmark.
\end{itemize}

If multiple benchmarks exist for the same task, researchers \should compare performance between benchmarks.
When selecting only a subset of all available benchmarks, researchers \should use the most specific benchmarks given the context.

Furthermore, researchers \should check whether a less resource-hungry approach (e.g., for static analysis tasks or program repair) can serve as a baseline. If so, the LLM or LLM-based tool should be compared to such baselines using suitable metrics. Even if LLM-based tools outperform baselines, researchers \should discuss whether the resources consumed justify the (often marginal) improvements~\cite{DBLP:journals/cacm/Menzies25}.


To compare traditional and LLM-based approaches or different LLM-based tools, researchers \should report established metrics whenever possible, as this allows secondary research.
They can, of course, report additional metrics that they consider appropriate.

In the following, we briefly discuss two common metrics used for generation tasks: \emph{BLEU-N} and \emph{pass@k}.
\emph{BLEU-N} \cite{DBLP:conf/acl/PapineniRWZ02} is a similarity score based on n-gram precision between two strings, ranging from $0$ to $1$.
Values close to $0$ represent dissimilar content, values closer to $1$ represent similar content, indicating that a model is more capable of generating the expected output.
%The score measures how many n-grams from the generated output of a model appear in the target string.
%In isolation, this approach favors shorter over longer generated sequences because it is more likely to have fewer n-grams in an extensive target sequence than more.
%To mitigate this, the so-called brevity penalty is introduced to disincentivize short sequences when having longer target sequences.
\emph{BLEU-N} has multiple variations.
\emph{CodeBLEU} \cite{DBLP:journals/corr/abs-2009-10297} and \emph{CrystalBLEU} \cite{DBLP:conf/kbse/EghbaliP22} are the most notable ones tailored to code.
They introduce additional heuristics such as AST matching.
As mentioned, researchers \must motivate why they chose a certain metric or variant thereof for their particular study.

The metric \emph{pass@k} reports the likelihood that a model correctly completes a code snippet at least once within \emph{k} attempts.
To our knowledge, the basic concept of \emph{pass@k} was first reported by \citeauthor{DBLP:journals/corr/abs-1906-04908} to evaluate code synthesis~\cite{DBLP:journals/corr/abs-1906-04908}.
They named the concept \emph{success rate at B}, where B denotes the trial ``budget.''
The term \emph{pass@k} was later popularized by \citeauthor{DBLP:journals/corr/abs-2107-03374} as a metric for code generation correctness~\cite{DBLP:journals/corr/abs-2107-03374}.
The exact definition of correctness varies depending on the task.
For code generation, correctness is often defined based on test cases: A passing test implies that the solution is correct.
The resulting pass rates range from $0$ to $1$.
A pass rate of $0$ indicates that the model was unable to generate a single correct solution within $k$ tries; a rate of $1$ indicates that the model successfully generated at least one correct solution in $k$ tries.

The \emph{pass@k} metric is defined as:

\begin{equation}
\text{pass@}k = 1 - \frac{\binom{n-c}{k}}{\binom{n}{k}}
\end{equation}

where $n$ is the total number of generated samples per prompt, $c$ is the number of correct samples among  $n$, and $k$ is the number of samples.

Choosing an appropriate value for \emph{k} depends on the downstream task of the model and how end-users interact with it.
% k= 1 -- pass@1
A high pass rate for \emph{pass@1} is highly desirable in tasks where the system presents only one solution or if a single solution requires high computational effort.
For example, code completion depends on a single prediction, since end users typically see only a single suggestion.
%In automated code refactoring, a large context might lead to high computational costs, making the calculation of a single solution expensive.
%Thus, having a high pass rate at pass@1 is crucial for such tasks.
% k = >=2 -- pass@{>=2}
Pass rates for higher $k$ values (e.g., $2$, $5$, $10$) indicate how well a model can solve a given task in multiple attempts.
%For downstream tasks that allow multiple solutions, e.g. via user interaction, strong performance at $k > 1$ can be justified. 
%For instance, a user selecting the correct test case from multiple suggestions allows for some model errors.
%In code search, finding the most appropriate good snippet for a given context, a reranking algorithm mitigate low pass rates at low k-values.
The \emph{pass@k} metric is frequently referenced in papers on code generation \cite{DBLP:journals/corr/abs-2308-12950, DBLP:journals/corr/abs-2401-14196, DBLP:journals/corr/abs-2409-12186, DBLP:journals/corr/abs-2305-06161}.
However, \emph{pass@k} is not a universal metric suitable for all generation tasks.
It requires a binary notion of correctness, making the metric appropriate for code synthesis evaluated via unit tests, but unsuitable for open-ended generation tasks such as comment generation, where multiple valid outputs exist.

If a study evaluates an LLM-based tool for supporting humans, a relevant metric is the acceptance rate, meaning the ratio of all accepted artifacts (e.g., test cases, code snippets) in relation to all artifacts that were generated and presented to the user.
Another way of evaluating LLM-based tools is calculating inter-model agreement (see also Section \href{/guidelines/#use-an-open-llm-as-a-baseline}{Use an Open LLM as a Baseline}).
This allows researchers to assess how dependent a tool's performance is on specific models and versions.
Metrics used to measure inter-model agreements include general agreement (percentage), \emph{Cohen's kappa}, and \emph{Scott's Pi coefficient}.

Due to LLM non-determinism, researchers \should repeat experiments to  assess (statistically) model or tool performance, e.g., using the arithmetic mean, median, confidence intervals, standard deviations, or more advanced statistical approaches~\cite{DBLP:conf/nips/AgarwalSCCB21, bjarnason2026randomnessagenticevals}.

From a measurement perspective, researchers \should reflect on the theories, values, and measurement models on which the benchmarks and metrics they have selected for their study are based.
For example, the phenomena labeled ``bugs'' in a large open dataset of software bugs is according to a certain theory of what constitutes a bug, and from the values and perspectives of the people who labeled the dataset. Reflecting on the context in which these labels were assigned and discussing whether and how the labels generalize to a new study context is crucial.

\guidelinesubsubsection{Example(s)}

According to \citeauthor{10.1145/3695988}, the main problem types for LLMs are classification, recommendation, and generation, each of which requires different metrics. In their systematic review,
\citeauthor{hu2025assessingadvancingbenchmarksevaluating} \cite{hu2025assessingadvancingbenchmarksevaluating} categorized 191 LLM benchmarks according to the specific SE task they address, making the paper a valuable resource for identifying existing metrics. 
Metrics used to assess generation tasks include \emph{BLEU}, \emph{pass@k}, \emph{Accuracy}, \emph{Accuracy@k}, and \emph{Exact Match}.
The most common metric for recommendation tasks is \emph{Mean Reciprocal Rank}.
For classification tasks, classical machine learning metrics such as \emph{Precision}, \emph{Recall}, \emph{F1-score}, and \emph{Accuracy} are often reported.

%\guidelinesubsubsection{Exact Match}
%\emph{Exact match} is another metric that calculates the percentage of replicas a model can produce.
%If the model is able to produce the exact target sequence, a score of 1 is awarded; otherwise 0.
%Compared to BLEU-N, \emph{exact match} is a stricter measurement.
%Due to its strictness, reported scores are usually low.

%he choice of metrics strongly depends on the specific task.
%Consequently, papers that address different research questions employ different metrics.
%For example, \citeauthor{DBLP:conf/nips/LiuXW023} assesses LLM models for their correctness in code generation using \emph{pass@k}~\cite{DBLP:conf/nips/LiuXW023}.
%However, \emph{pass@k} is not an universal metric suitable for every problem at hand.
%For instance, for creative or open-ended tasks such as comment generation, pass@k is not a suitable metric since not only one single correct result exists.
%Thus, \citeauthor{DBLP:conf/icse/GengWD00JML24} evaluates LLMs for code comment generation using BLEU, ROUGE-L and METEOR as evaluation %metrics~\cite{DBLP:conf/icse/GengWD00JML24}.
%\citeauthor{DBLP:journals/ase/LiuJZNLL25} investigates the performance of LLMs on automated refactorings tasks and proposes a novel metric tailored to the unique characteristics of this problem \cite{DBLP:journals/ase/LiuJZNLL25}.

%Code snippets provided by benchmarks contain untrusted code, making sandboxing recommendable.
%To avoid potential security risks, execute untrusted code in an isolated environment, such as a virtual machine or containerization, instead of your own environment.

Benchmarks used for code generation incldue \emph{HumanEval} (available on \href{https://github.com/openai/human-eval}{GitHub}) \cite{DBLP:journals/corr/abs-2107-03374}, \emph{MBPP} (available on \href{https://huggingface.co/datasets/google-research-datasets/mbpp}{Hugging Face}) \cite{DBLP:journals/corr/abs-2108-07732},
% Both benchmarks consist of code snippets written in Python sourced from publicly available repositories.
% Each snippet consists of four parts: a prompt containing a function definition and a corresponding description of what the function should accomplish, a canonical solution, an entry point for execution, and test cases.
% The input of the LLM is the entire prompt.
% The output of the LLM is then evaluated against the canonical solution using metrics or against a test suite.
\emph{ClassEval} (available on \href{https://github.com/FudanSELab/ClassEval}{GitHub}) \cite{DBLP:journals/corr/abs-2308-01861}, \emph{LiveCodeBench} (available on \href{https://github.com/LiveCodeBench/LiveCodeBench}{GitHub}) \cite{DBLP:journals/corr/abs-2403-07974}, and \emph{SWE-bench} (available on \href{https://github.com/swe-bench/SWE-bench}{GitHub}) \cite{DBLP:conf/iclr/JimenezYWYPPN24}.
An example of a code translation benchmark is \emph{TransCoder}~\cite{DBLP:journals/corr/abs-2006-03511} (available on \href{https://github.com/facebookresearch/CodeGen}{GitHub}). 

% Lukas Boehme: Add examples for other categories: classification and recommendation
Given the generative nature of LLMs, most benchmarks focus on generation tasks, but benchmarks for classification and recommendation tasks also exist.
For classification tasks such as vulnerability detection, \emph{DiverseVul} (available on \href{https://github.com/wagner-group/diversevul}{GitHub})~\cite{DBLP:conf/raid/0001DACW23} provides a dataset of vulnerable and non-vulnerable functions.
They evaluate the binary classification using  \emph{Accuracy}, \emph{Precision}, \emph{Recall}, \emph{F1-score} and \emph{False Positive Rate}.
For recommendation tasks such as code search, \emph{CodeSearchNet} (available on \href{https://github.com/github/CodeSearchNet}{GitHub})~\cite{DBLP:journals/corr/abs-1909-09436} contains code-documentation pairs across multiple programming languages.
The recommendations are evaluated using \emph{Mean Reciprocal Rank}.

\guidelinesubsubsection{Benefits}

Benchmarks and metrics improve reproducibility and comparability across studies, leading to faster improvements in the cumulative body of knowledge. It is simply easier to assess a new system against one or a few benchmarks than hundreds of competing systems, and when most systems are assessed against the same benchmarks and metrics, we can see which approaches work best (at least, on those benchmarks). This comparison enables progress tracking; for example, when researchers iteratively improve a new LLM-based tool and test it against benchmarks after significant changes. For practitioners, leaderboards (published benchmark results of models) support selecting models for downstream tasks. \added{Meanwhile, baselines help us understand just how much LLMs improve performance over non-LLM-enabled alternatives.} 

\guidelinesubsubsection{Challenges}

\added{Computer scientists have a long history of inventing oversimplified, unvalidated metrics (e.g. lines of code, CPU time) and simply claiming that they are reasonable proxies for complicated, value laden, under-theorized, often multidimensional, latent variables (e.g. system size, environmental sustainability). In scientific fields that take measurement more seriously, researchers expect that metrics and instruments are backed up by foundational theories (e.g. we accept telescopes based on our understanding of the physics of light) or empirical evidence (e.g. we accept questionnaire scales for the Big Five personality traits based on extensive quantitative and qualitative investigations of their psychometric properties). The fundamental challenge with AI metrics and benchmarks is that many of them have neither firmly understood theoretical underpinnings nor extensive empirical validation of their construct, measurement, and ecological validity.} 

For example, metrics such as \emph{pass@k} do not reflect qualitative aspects of the generated source code, including its maintainability or readability.
However, these aspects are critical for many downstream tasks.
Moreover, benchmarks are isolated test sets and may not fully represent real-world applications.
For example, benchmarks such as \emph{HumanEval} synthesize code based on written specifications.
However, such explicit descriptions are rare in real-world applications.
Thus, evaluating model performance with benchmarks might not reflect real-world tasks and end-user performance.

Benchmarks and metrics also have generalizability problems. Prominent LLM benchmarks such as \emph{HumanEval} and \emph{MBPP} use Python, so researchers can optimize for Python's idiosyncrasies, creating illusory improvements that do not generalize to other languages. Similarly, the metric \emph{BLEU-N} is a syntactic metric, so code can score highly without being executable. The metric \emph{Exact Match}, meanwhile, does not account for functional equivalence of syntactically different code. Both \emph{BLEU-N} and \emph{Exact Match} are influenced by code formatting, which confounds their intended use. 

Execution-based metrics such as \emph{pass@k} directly evaluate correctness by running test cases, but they require a setup with an execution environment.
When researchers observe unexpected values for certain metrics, the specific results should be investigated in more detail to uncover potential problems.


% Many closed-source models, such as those released by \emph{OpenAI}, achieve exceptional performance on certain tasks but lack transparency and reproducibility~\cite{DBLP:conf/nips/00110ZZDJLHL24, DBLP:journals/corr/abs-2308-01861, DBLP:journals/corr/abs-2406-15877}.
% Benchmark leaderboards, particularly for code generation, are led by closed-source models~\cite{DBLP:journals/corr/abs-2308-01861, DBLP:journals/corr/abs-2406-15877}.
% While researchers \should compare performance against these models, they must consider that providers might discontinue them or apply undisclosed pre- or post-processing beyond the researchers' control (see \openllm).
%The lack of transparency introduces challenges in analyzing mistakes, limiting the understanding of model failures.
%Open-source models offer greater transparency since the entire process is under the researcher's control; however, they require appropriate hardware.

\added{Finally, benchmark data contamination---where the benchmark itself is part of the training data---may lead to artificially high performance if the model remembers the solution from the training data rather than actually solving the task based on the input data~\cite{DBLP:journals/corr/abs-2406-04244}. Therefore, for proprietary LLMs that do not release their training data, researchers \should consider using human validation (Section~\ref{sec:use-human-validation-for-llm-outputs}), curating new data, or refactoring existing data, or~\cite{DBLP:journals/corr/abs-2406-04244}.} 

\added{Creating and curating new uncontaminated benchmark datasets produces unseen benchmark data.
This can be achieved by collecting data after a specified cutoff date, ensuring that models trained before this date could not have been exposed to the benchmark. To ensure transparency, researchers \must disclose the data sources and their collection dates for each benchmark release. However, this temporal approach may be impractical because it requires continuous updates as new models might include the benchmark in future training datasets. Keeping the benchmark private prevents inclusions in training sets, yet it requires trust in the benchmark creator and a system to execute the benchmark without leaking the benchmark.}

\added{Refactoring existing data restructures benchmarks by generating new prompts or augmenting samples across multiple dimensions. Thus, pure memorization by the model is ineffective. However, this only partially addresses contamination since the model may have learned underlying concepts from other versions of the benchmark.}


%\removed{Models may become overfitted to the training data, leading to poor performance on unseen data. When optimizing for a metric, the underlying model might adjust its parameters to improve the metric, thereby failing to generalize. This can result in a model performing exceptionally well on a benchmark but struggling to hold up the performance in unforeseen settings, such as real-world scenarios.}

\guidelinesubsubsection{Study Types}

This guideline \must be followed for all study types that automatically evaluate the performance of LLMs or LLM-based tools.
The design of a benchmark and the selection of appropriate metrics are highly dependent on the specific study type and research goal.
Recommending specific metrics for specific study types is beyond the scope of these guidelines, but \citeauthor{hu2025assessingadvancingbenchmarksevaluating} provide a good overview of existing metrics for evaluating LLMs~\cite{hu2025assessingadvancingbenchmarksevaluating}.
In addition, our Section \humanvalidation provides an overview of the integration of human evaluation in LLM-based studies. 
For \annotators, the research goal might be to assess which model comes close to a ground truth dataset created by human annotators.
Especially for open annotation tasks, selecting suitable metrics to compare LLM-generated and human-generated labels is important.
In general, annotation tasks can vary significantly.
Are multiple labels allowed for the same sequence? Are the available labels predefined, or should the LLM generate a set of labels independently?
Due to this task dependence, researchers \must justify their metric choice, explaining what aspects of the task it captures together with known limitations.
If researchers assess a well-established task such as code generation, they \should report standard metrics such as \emph{pass@k} and compare the performance between models.
If non-standard metrics are used, researchers \must state their reasoning.

\guidelinesubsubsection{Advice for Reviewers}

\added{Reviewers should expect manuscripts to: 
\begin{enumerate}
    \item clearly identify the constructs or variables that the study ought to measure in principle (e.g. LLM performance, quality of generated code), including independent variable(s), dependent variable(s) and (both physical and statistical) control variables;
    \item present their measurement model; that is, identify exactly which metrics, benchmarks, or baselines are used and how they relate to constructs or variables the study aims to measure;
    \item explicitly justify the selection of any metrics, benchmarks, and baselines used
    item contemplate \textit{in detail} the assumptions, reliability, and validity---\textit{especially construct validity} of each benchmark and metric;
    \item clearly articulate any limitations regard construct and measurement validity (see also Section~\ref{sec:report-limitations-and-mitigations}).
\end{enumerate}}

\added{As in previous guidelines, if some specific information about a baseline, etc. is missing, reviewers should simply ask authors to add it. In contrast, reviewers must not tolerate vague descriptions that conflate what researchers and practitioners actually care about (e.g. effectiveness, efficiency, quality, etc.) with specific methods of counting something ostensibly indicative of these broader concerns. Reviewers should be wary of popularity arguments; while the whole point of benchmarks is that standardizing evaluation is beneficial, ubiquity does not imply validity or appropriateness for any given context. Similarly, reviewers must not tolerate ambivalence toward construct validity or obfuscation of validity concerns.  In summary, for this guideline, reviewers should look for manuscripts to convey a solid understanding of construct and measurement validity by explaining and justifying selected measurement models. }